% Nouveau EGA V (cote 25): archive cote 25, site page 158 (PDF page 159).
% typescript. Draft compared: none.
% AI-transcribed by Claude Opus 5.5 (claude-opus-5-5, adaptive thinking, effort high), Message Batches API, batch ega25b (msgbatch_01A7n5YuXubUPNK1zsepi8RG), 2026-11-13; not yet checked by eye.
\scanpage{25/158}
\margin{nouveau EGA V} \mspage{1}

IV 17.15. \uline{Formes lisses}; \inserted{\unsure{augmentations|argumentations} et singularités élémentaires.} \deleted{\xout{x} formes quadratiques lisses}

NB Je viens de m'apercevoir que la terminologie introduite dans
ma rédaction \inserted[typed]{SUP} de la Supersingularité est déraisonnable, et conflicte
en particulier avec la terminologie recue. (De toutes façons, tu a\unsure{\xout{s}} dû
no\over{t}{t}er que dans la déf.1 de ce texte, il faut lire ``dégénérée'' au lieu de
``non dégénérée''). Le canular provient des variétés de dimension paire
en car.2 , \xout{\unsure{une telle variété}} définie\hand{s} par exemple par une équation $\phi=0$
dans une variété ambiante de dim impaire, car avec la terminologie de mes
notes, une telle variété ne peut avoir de \xout{\unsure{zéro sing}} \hand{``}point singulier
or\over{d}{d}inaire\hand{''} (= singularité quadratique ordinaire), i.e. $\phi$ ne peut avoir
de \hand{``}zéro singulier ordinaire\hand{''}, d\over{u}{û} au fait qu'en car 2, une forme quadratique
à un nombre impair de variables est toujours ``dégénérée''. Or tout le
monde a toujours considéré que même en car. 2, l'origine du cône affine
$x^2 + yz = 0$ est une singularité quadratique ordinaire. Dans les présen-%
tes notes, je dégage la notion de \xout{f}forme quadratique lisse (ou ``ordinaire''),
(sur un Module loc libre de type fini), de telle façon que
non dégénéré \hand{$\Longrightarrow$} lisse, la réciproque étant vraie si E est de rang pair
\uline{ou} les car résiduelles de la base S sont toutes $\neq 2$ . Cette notion étant
posée, je propose de garder la terminologie ``supersingulière'' (qui ne
conflicte avec aucune terminologie reçue) \xout{\unsure{xxxxxx de mon texte}} de SUP,
comme correspondant à la notion de forme quadratique non dégéné\over{f}{r}ée; on
parlera de même de point supersingulier d'un préschéma loc noeth. , ou
géométriquement supersingulier pour un préschéma loc de type fini sur un
corps, dans le même sens (en regardant l'hom $\mathrm{Sym}(\underline{M}/\underline{m}^2) \hand{\longrightarrow} \mathrm{gr}_{\underline{m}}(A)$ et
\unsure{en|ex} disant que le noyau ne peut être engendré par une forme quadratique
non dégénérée); donc ce qui dans SUP s'appelait à tort ``zéro singulier
ordinair\hand{e}'' s'appelera ``zéro singulier non supersingulier'' , ou ``zéro

\margin{158}
